\section{Annihilator-aware section capacity}\label{sec:capacity}
All sections in this section are sections of one specified permutation
module. This restriction supplies the capacity bound; odd order alone
does not supply it.

\subsection{Permutation-module bounds}
For a transitive group of degree $u$, write $u_2$ for the binary part
of $u$, and set $\lpp(1)=1$. The induced-module theorem of
\cite{Tracey2018} bounds the number of generators of any submodule of
the binary permutation module by
\[
 E(u,2)=\min\left\{
 \left\lfloor\sqrt{2/\pi}\frac{u}{\sqrt{\log u_2}}\right\rfloor,
 \frac{u}{\lpp(u/u_2)}\right\}
\]
when no soluble transitive subgroup exists. The first entry is infinite
if $u_2=1$. When such a subgroup exists the applicable bound is
\[
 E_{\rm sol}(u,2)=\min\left\{
 \frac{u}{2^{K(u)}}\binom{K(u)}{\lfloor K(u)/2\rfloor},u_2\right\},
 \qquad K(u)=\sum_{p^a\parallel u}a(p-1).
\]
Both formulas imply the common integer bound $d_T(L)\leq E(u,2)$:
use $K(u)\geq\log u_2$, the central-binomial estimate,
$u_2\leq u/\lpp(u/u_2)$, and integrality of generator number.
The sharper soluble formula will be needed as well.

Directly from these formulas,
\begin{equation}\label{eq:E-consequences}
 E(u,2)\leq u/2\ (u\geq2),\qquad
 E(u,2)\leq3u/8\ (u\geq8),\qquad
 E(u,2)\leq11u/32\ (u\geq24\text{ even}).
\end{equation}
For a nontrivial odd part use $E\leq u/3$. For powers of two the
degrees $2,4,8,16,32$ give respectively $1,2,3,6,11$; from $64$
onward the square-root expression is less than $u/3$.

\begin{lemma}[Orbit filtration]\label{lem:orbit-filtration}
Let $H\lhd T$ have $k$ orbits of length $u\geq2$ in a transitive
action of degree $s=ku$. Every trivial $H$-section of $\F_2^s$ has
dimension at most $kE(u,2)$.
\end{lemma}
\begin{proof}
Filter a submodule by its intersections with successive sums of the
actual orbit modules. Each successive section embeds in one orbit
module. Right exactness of coinvariants bounds the dimension of the
coinvariants of the original submodule by the sum of the corresponding
dimensions of these sections. Each is bounded by its module-generator
number and hence by $E(u,2)$. A trivial quotient factors through
coinvariants. This filtration does not require independent orbit images.
\end{proof}

\begin{lemma}[A sharper fixed-section bound]\label{lem:five-sixteenths}
If $T$ is transitive of even degree $s\geq24$, every submodule $L$ of
$\F_2^s$ satisfies $d_T(L)\leq5s/16$. Consequently every section $A$
satisfies $\dim A^T\leq\lfloor5s/16\rfloor$.
\end{lemma}
\begin{proof}
Write $s=2^ao$ with $o$ odd. If $o\notin\{1,3\}$, then
$\lpp(o)\geq5$, and the two appropriate bounds are respectively
$E\leq s/5$ and $E_{\rm sol}\leq s/o\leq s/5$.
If $o=1$, a Sylow $2$-subgroup is transitive: for a point $x$ its
stabilizer in a Sylow subgroup has order at most $|T_x|_2=|T|_2/s$.
Here $a\geq5$. If $o=3$ and a soluble transitive subgroup exists, then
$a\geq3$ and $K(s)=a+2\geq5$. In both cases the central-binomial
proportion is at most $\binom52/2^5=5/16$; it is nonincreasing.

It remains to consider $s=3\cdot2^a$ without a soluble transitive
subgroup. Apply \cite[Corollary 3.12]{Tracey2018} to an actual minimally
transitive subgroup. It supplies a soluble subgroup with, for
$0\leq j\leq e$, exactly $\binom ej2^{u_1}$ orbits of length
$3p^j2^{u-u_1}$, where $p=2^h-1$ is a Mersenne prime,
$h\geq3$, $e\geq1$, $a=eh+u$, and $u\geq u_1\geq0$.
Filter by these actual orbit modules. The soluble bound on each is at
most its binary part, giving
\[
 d_T(L)\leq2^{e+u_1}2^{u-u_1}=2^{e+u}\leq s/12.
\]
Restriction can only increase the required generator number, and
generator number is subadditive on the filtration. Finally lift $A^T$
to a submodule of $\F_2^s$; its trivial quotient requires
$\dim A^T$ generators.
\end{proof}

\begin{lemma}[Unipotent sections]\label{lem:unipotent}
Suppose that $T\leq S_s$ is faithful and transitive and is not a
$2$-group. If every composition factor of a section $P$ of $\F_2^s$
is trivial, then $\dim P\leq3s/8$.
\end{lemma}
\begin{proof}
The action on $P$ is conjugate to a group of unitriangular matrices,
so $O^2(T)$ acts trivially. This is a nontrivial normal subgroup.
It is generated by odd-order elements, as are its images on its
orbits; each nontrivial transitive orbit image is therefore nonbinary.
All orbit lengths are equal, say $u$.
For $u\geq8$, use Lemma~\ref{lem:orbit-filtration}.
For $u=3,5,6,7$, choose an odd prime dividing $u$. Its Sylow subgroup
has no fixed points, since the $p$-part of a point stabilizer is
strictly smaller than the $p$-part of the group. Exactness of invariants
for this odd group bounds any trivial section by the number of its
orbits, at most $u/p\leq3u/8$.

For $u=4$, a nonbinary transitive subgroup contains a $3$-cycle with
fixed point $x$. Its invariant subspace has dimension two. A trivial
section of dimension two would force the numerator to contain this
entire invariant subspace, including the coordinate vector at $x$.
Transitivity would force the numerator to be the whole permutation
module, whose largest trivial quotient has dimension one. This is a
contradiction. Thus the bound holds at $u=4$ too. Sum over the actual
orbit filtration as in Lemma~\ref{lem:orbit-filtration}.
\end{proof}

\subsection{A coupled socle estimate}
\begin{proposition}\label{prop:coupled}
For every section $M$ of the binary permutation module of a faithful
transitive group $T$ of even degree $s\geq24$, put
$t=\dim M^T$ and
$r_{\rm nt}(M)=\max_{X\ne1}m_X(\Soc M)/\dim_{\End_TX}X$.
Then
\begin{equation}\label{eq:coupled}
 2t+4r_{\rm nt}(M)\leq47s/48.
\end{equation}
No splitting condition on $M^T$ is required.
\end{proposition}
\begin{proof}
We have $t\leq11s/32$ by \eqref{eq:E-consequences}. If there is a
nontrivial simple type, choose $X$ attaining the maximum and write
$h=\dim_{\F_2}X$, $\End_TX=\F_{2^f}$, $d=h/f$, $m=m_X$,
and $H=\ker(T\to\GL(X))$. The mixed section $1^t\oplus X^m$
is the literal sum in $M$ of $M^T$ and the full $X$-isotypic socle
component.  The sum is direct: a nonzero fixed vector in the latter
would generate a trivial simple submodule and hence force $X=1$.
Thus it has dimension $L=t+hm$; no complement to $M^T$ is being chosen.

If $H\ne1$, it has orbit length at least two, and
Lemma~\ref{lem:orbit-filtration} gives $L\leq s/2$.
For $hd\geq3$,
\[
 2t+4m/d=\frac4{hd}L+\left(2-\frac4{hd}\right)t
 \leq\frac{11s}{16}+\frac{5s}{8hd}\leq\frac{43s}{48}.
\]
The only remaining nontrivial possibility is $h=2,d=1$.
The image is then $C_3$, acting as nonzero scalars of $\F_4$.
There are one or three $H$-orbits, each of length at least eight.
Thus $L\leq3s/8$, and the required cost is $2L\leq3s/4$.

If $H=1$, the group embeds in $\GL_d(2^f)$. Since its degree is
even, $d\ne1$. The values $hd=fd^2<9$ are $4$ and $8$.
The first would give $|T|\leq6$, impossible. In the second case
$T\leq\GL_2(4)$, of order $180$. Since $s\mid |T|$ and $s\geq24$,
the odd part of $s$ has largest prime-power divisor at least five.
Thus $t\leq s/5$, and $L\leq s$ gives
$2t+4m/d\leq s/2+3t/2\leq4s/5$.
For $hd\geq9$, the same dimension bound gives
\[
 2t+4m/d\leq\frac{11s}{16}+\frac{21s}{8hd}
 \leq\frac{47s}{48}.
\]
If there is no nontrivial type, the initial bound on $t$ suffices.
\end{proof}

\subsection{The retained annihilator and the central cut}
For a finite $\F_2B$-module $A$, put $S=A^B$ and, for each simple
$X$, define the actual module-extension annihilator
\[
 \mathcal L_X=\ker\bigl(\operatorname{Ext}^1_B(X,S)
                  \longrightarrow\operatorname{Ext}^1_B(X,A)\bigr).
\]
Every $C\leq S$ has a module complement in $S$, so
$\operatorname{Ext}^1_B(X,C)$ embeds as a specified subspace of
$\operatorname{Ext}^1_B(X,S)$. The long exact sequence gives
\begin{equation}\label{eq:exact-socle}
 m_X(A/C)=m_X(A)-\mathbf1_{X=1}\dim C+
 \dim_{\End_BX}\bigl(\mathcal L_X\cap\operatorname{Ext}^1_B(X,C)\bigr).
\end{equation}
Here multiplicity always means socle multiplicity. In particular
\begin{equation}\label{eq:fixed-quotient}
 \dim(A/C)^B=\dim A^B-\dim C+
 \dim\ker\bigl(H^1(B,C)\to H^1(B,A)\bigr).
\end{equation}
These identities explicitly account for newly exposed simple constituents.

\begin{proof}[Proof of Theorem~\ref{thm:capacity-intro}]
Let $t=\dim S$, $\mathcal L_1=\ker(H^1(B,S)\to H^1(B,A))$,
and $\ell=\dim\mathcal L_1$. The invariants sequence identifies
$\mathcal L_1$ with $(A/S)^B$. Its preimage in $A$ is a section of
dimension $t+\ell$ with only trivial composition factors. Hence
\begin{equation}\label{eq:two-budgets}
 t\leq\lfloor5s/16\rfloor,\qquad t+\ell\leq3s/8
\end{equation}
by Lemmas~\ref{lem:five-sixteenths} and~\ref{lem:unipotent}.

Use the natural identification
$H^1(B,\F_2)\otimes S\cong
\Hom(S^*,H^1(B,\F_2))$ and regard $\mathcal L_1$ as an
$\ell$-dimensional space of maps out of $S^*$.  The elementary
evaluation-separator lemma gives a subspace $W\leq S^*$ with
$\dim W\leq\ell$ such that restriction to $W$ is injective on
$\mathcal L_1$: starting with the whole map space, choose an input on
which a nonzero remaining map does not vanish and iterate on the kernel
of that evaluation.  Put $Z=W^\perp\leq S$. Then
\[
 \dim Z\geq\max(t-\ell,0),\qquad
 \mathcal L_1\cap H^1(B,Z)=0.
\]
Indeed, a tensor in $H^1(B,Z)$, viewed as a map on $S^*$, vanishes on
$Z^\perp=W$, so injectivity forces every element of the displayed
intersection to be zero. Choose $C\leq Z$ of dimension
\[
 c=\min\{\lfloor t/2\rfloor,\max(t-\ell,0)\}.
\]
Equation~\eqref{eq:fixed-quotient} gives $\dim(A/C)^B=t-c$.
Apply Proposition~\ref{prop:coupled} to this actual quotient:
\[
 2c+4r_{\rm nt}(A/C)\leq47s/48-2t+4c\leq47s/48.
\]
For its trivial type the cost is
\[
 2c+4(t-c)=4t-2c\leq\max\{3t+1,2(t+\ell)\}.
\]
The inequality follows by considering which entry attains the minimum
defining $c$, including $\ell\geq t$. By \eqref{eq:two-budgets},
\[
 3t+1\leq15s/16+1\leq47s/48,\qquad
 2(t+\ell)\leq3s/4.
\]
The first inequality uses $s\geq24$. Taking the maximum over the
trivial and nontrivial types proves \eqref{eq:capacity-intro} with one
and the same $C$.
\end{proof}

\begin{remark}[Why the permutation section matters]
On $\F_2^t\oplus\F_2^t$ with trivial $C_3$ action, every linear graph
from one summand to the other is invariant, giving $2^{t^2}$ graphs.
Thus odd order supplies no strict quadratic deficit in an arbitrary
trivial representation. The proof above instead uses
\eqref{eq:two-budgets}, the actual annihilator, and a quotient cut.
It does not assert a strict deficit for the ranks of the uncut trivial
axes.
\end{remark}
